% \subsection{Đóng góp dự kiến}

% Nghiên cứu này dự kiến mang lại các đóng góp sau:
% \begin{enumerate}
%     \item Thực nghiệm hệ thống mô hình DeepLabV3 trên tập dữ liệu FoodSeg103.
%     \item Đánh giá hiệu quả của các backbone lightweight trong bài toán phân đoạn thực phẩm.
%     \item Phân tích sự cân bằng giữa độ chính xác và tốc độ suy luận.
%     \item Đề xuất cấu hình mô hình phù hợp cho triển khai trên thiết bị di động.
% \end{enumerate}

% \subsection{Tài nguyên cần thiết}

% \textbf{Dữ liệu:} FoodSeg103 dataset.

% \textbf{Phần cứng:} GPU 16GB VRAM (Colab T4/Kaggle P100).

% \textbf{Phần mềm:} PyTorch, ExecuTorch.


\section*{Đóng góp dự kiến}

Nghiên cứu kỳ vọng mang lại ba đóng góp chính. Thứ nhất, cung cấp một đánh giá có hệ thống về mức độ phù hợp của các mô hình phân đoạn ngữ nghĩa CNN-based nhẹ đối với bài toán phân đoạn ảnh món ăn fine-grained trên FoodSeg103. Thứ hai, làm rõ hơn mối quan hệ giữa hiệu quả tính toán và khả năng biểu diễn đặc trưng texture trong các mô hình lightweight. Thứ ba, đề xuất một cấu hình hoặc hướng cải tiến theo texture-aware có tính khả thi cho các kịch bản triển khai thực tế trên thiết bị có tài nguyên hạn chế.

\section*{Tài nguyên cần thiết}
\textbf{Dữ liệu:} FoodSeg103.  

\textbf{Phần cứng:} GPU có bộ nhớ từ 16GB VRAM trở lên.  

\textbf{Phần mềm:} PyTorch và các công cụ phục vụ đánh giá, tối ưu suy luận và triển khai mô hình.